\section*{The Propers: Detailed Commentary}

\subsection*{Gathered before being corrected}

The Entrance's scriptural adaptation places the assembly in God's holy place
and under divine strength. This matters for the narrative that follows:
Israel's murmuring and the crowd's misdirected search occur within a history
already sustained by gift. The Collect likewise asks for protection,
governance, and guidance before asking for right action. Its passing/enduring
contrast does not license contempt for created goods; it asks that they be
used without losing the goods that last.

\subsection*{Exodus and Psalm: manna remembered truthfully}

Exodus 16 holds complaint, provision, and testing together. The people
remember Egyptian food while suppressing bondage; God nevertheless hears and
gives. The daily sufficiency of v.~4 prevents hoarding from defining the gift,
while the people's question in v.~15 preserves its strangeness. Psalm 78
receives the episode as intergenerational instruction. Augustine's exposition
of Psalm 77 distinguishes the historical benefits from the spiritual reality
to which they point, while also treating forgetfulness as the psalm's moral
danger. The response therefore makes memory itself an act of worship.

\subsection*{Ephesians: the form of renewed reception}

The omitted vv.~18--19 must not be silently supplied as though proclaimed.
Within the letter, however, the selected boundary belongs to a larger contrast
between gentile futility and life learned in Christ. ``Put off'' and ``put
on'' are not cosmetic metaphors: Chrysostom's \emph{Homily on Ephesians} 13
treats renewal of the mind as a transformation expressed in conduct. The
newness is created ``according to God,'' so renewal is received before it is
performed; it nonetheless issues in justice and holiness of truth.

\subsection*{John 6: seeking, working, believing}

John's scene follows the feeding and crossing. In Homilies 43--45 on John,
Chrysostom tracks the passage in sequence: vv.~24--25 show the crowd's pursuit;
vv.~26--27 expose satisfaction with loaves and redirect labor toward enduring
food; vv.~28--35 answer the demand for work and signs by centering belief in
the sent Son and the Father's gift. His moral insistence on purified motive
should not be converted into disdain for the hungry. Jesus himself has fed
the crowd; the correction concerns what they perceive in the sign.

Augustine's \emph{Tractate on John} 25, especially sections 8--14, likewise
distinguishes seeking Christ for temporal advantage from coming to him in
faith. His reading of eating as believing clarifies the discourse's movement
without exhausting its sacramental reception. The crowd cites manna as a
demand for another sign; Jesus shifts the grammar from ancestral past to the
Father's present giving. Verse 35 then identifies bread with the speaker:
coming and believing parallel hunger and thirst.

\subsection*{Acclamation and eucharistic conclusion}

Matthew 4:4b is an acclamatory adaptation from the temptation narrative. It
supports the distinction between bread alone and life from God's word, but
does not relocate John's crowd to Matthew's wilderness. At the altar, the
offerings prayer confesses divine bounty before human return. Ambrose's
\emph{On the Mysteries} 8.44--49 illuminates sacramental reception without
directly exegeting every appointed verse. Communion B repeats the Gospel's
climax, while Communion A turns to grateful memory. The final prayer names
the Passion memorial and asks for fruitful reception, joining remembrance
and transformation.
